% Part: second-order-logic
% Chapter: syntax-and-semantics
% Section: expressive-power

\documentclass[../../../include/open-logic-section]{subfiles}

\begin{document}

\olfileid{sol}{syn}{exp}
\olsection{القوة التعبيرية}

\begin{explain}
التكميم على متغيرات الرتبة الثانية مسؤول عن زيادة هائلة في القوة
التعبيرية للغة مقارنة بمنطق الرتبة الأولى. فالتكميم الوجودي من الرتبة
الثانية يتيح لنا أن نقول إن دوال أو علاقات لها خصائص معينة موجودة. ولا
سبيل إلى ذلك في منطق الرتبة الأولى إلا بتخصيص رمز غير منطقي (أي
!!a{function} أو~!!{predicate}) لهذا الغرض. ويتيح لنا التكميم الكلي
من الرتبة الثانية أن نقول إن جميع المجموعات الجزئية من الـ!!{domain}،
أو العلاقات عليه، أو الدوال منه إلى الـ!!{domain} لها خاصية ما. أما في
منطق الرتبة الأولى فلا يمكننا إلا القول إن للمجموعات الجزئية أو
العلاقات أو الدوال المسندة إلى أحد الرموز غير المنطقية للغة خاصية ما.
وحين نقول إن مجموعات جزئية أو علاقات أو دوال لها خاصية ما موجودة، أو
إنها كلها تملك تلك الخاصية، يمكننا أن نستعمل التكميم من الرتبة الثانية
في تحديد الخاصية نفسها أيضًا. وهذا يتيح لنا تعريف علاقات لا تقبل
التعريف في منطق الرتبة الأولى، والتعبير عن خصائص للمجال لا يمكن التعبير
عنها فيه.
\end{explain}

\begin{defn}
إذا كانت~$\Struct{M}$ !!a{structure} للغة~$\Lang{L}$، كانت العلاقة
$R \subseteq \Domain{M}^2$ \emph{قابلة للتعريف} في~$\Lang{L}$ إذا
وجدت~!!{formula}~$!A_R(x, y)$ لا تحوي من المتغيرات الحرة إلا~$x$
و$y$، بحيث تصدق~$R(a, b)$ (أي~$\tuple{a,b} \in R$) إذا وفقط إذا
$\Sat{M}{!A_R(x, y)}[s]$ عندما~$s(x) = a$ و$s(y) = b$.
\end{defn}

\begin{ex}
يمكننا في منطق الرتبة الأولى تعريف علاقة الهوية
$\Id{\Domain{M}}$ (أي~$\Setabs{\tuple{a,a}}{a \in
  \Domain{M}}$) بالصيغة~$\eq[x][y]$. ويمكننا في منطق الرتبة الثانية
تعريف هذه العلاقة \emph{من دون~$\eq$}. فإذا كان~$a$ و$b$ هما
الـ!!{element} نفسه من~$\Domain{M}$، كانا~!!{element}s في المجموعات
الجزئية نفسها من~$\Domain{M}$ (لأن المجموعات تتعين بـ!!{element}s
فيها). وبالعكس، إذا اختلف~$a$ و$b$، لم يكونا~!!{element}s في
المجموعات الجزئية نفسها: فمثلًا~$a \in \{a\}$ لكن
$b \notin \{a\}$ إذا كان~$a \neq b$. ولذلك فإن «كونهما
!!{element}s في المجموعات الجزئية نفسها من~$\Domain{M}$» علاقة تصدق
على~$a$ و$b$ إذا وفقط إذا كان~$a = b$. ويمكن التعبير عن هذه العلاقة
في منطق الرتبة الثانية، إذ نستطيع التكميم على جميع المجموعات الجزئية
من~$\Domain{M}$. ومن ثم تعرّف الـ!!{formula} الآتية
$\Id{\Domain{M}}$:
\[
\lforall[X][(X(x) \liff X(y))]
\]
\end{ex}

\begin{prob}
بيّن أن~$\lforall[X][(X(x) \lif X(y))]$ (لاحظ: $\lif$ لا~$\liff$!)
تعرّف~$\Id{\Domain{M}}$.
\end{prob}

\begin{ex}
إذا كان~$R$ !!{predicate} ثنائي الموضع، كانت~$\Assign{R}{M}$ علاقة
ثنائية الموضع على~$\Domain{M}$. وربما على نحو يوقع في شيء من الالتباس،
سنستعمل~$R$ لكل من الـ!!{predicate}~$R$ والعلاقة~$\Assign{R}{M}$
نفسها. ويكون \emph{الإغلاق المتعدي}~$R^*$ للعلاقة~$R$ هو العلاقة
التي تصدق بين~$a$ و$b$ إذا وفقط إذا وجدت
$c_1$، و\dots، و$c_k$ بحيث تصدق
$R(a,c_1)$، و$R(c_1, c_2)$، و\dots، و$R(c_k,b)$. ويشمل ذلك حالة
$k = 0$؛ أي إنه إذا صدقت~$R(a,b)$ صدقت~$R^*(a,b)$ أيضًا. ومن ثم
$R \subseteq R^*$. والواقع أن~$R^*$ أصغر علاقة تحتوي~$R$ وتكون
متعدية. ويمكننا أن نقول في منطق الرتبة الثانية إن~$X$ علاقة متعدية
تحتوي~$R$:
\begin{multline*}
  !B_R(X) \ident \lforall[x][\lforall[y][(R(x,y) \lif X(x, y))]] \land {}\\
\lforall[x][\lforall[y][\lforall[z][((X(x,y) \land X(y,z)) \lif X(x,
      z))]]].
\end{multline*}
يقول العامل الأول إن~$R \subseteq X$، ويقول الثاني إن~$X$ متعدية.

وللقول إن~$X$ أصغر علاقة من هذا القبيل، نقول إنها نفسها محتواة في كل
علاقة تحتوي~$R$ وتكون متعدية. ولذلك يمكننا تعريف الإغلاق المتعدي
لـ$R$ بالـ!!{formula}
\[
R^*(X) \ident !B_R(X) \land \lforall[Y][(!B_R(Y) \lif
  \lforall[x][\lforall[y][(X(x, y) \lif Y(x,y))]])].
\]
لدينا~$\Sat{M}{R^*(X)}[s]$ إذا وفقط إذا~$s(X) = R^*$. ولا يمكن
التعبير عن الإغلاق المتعدي لـ$R$ في منطق الرتبة الأولى.
\end{ex}

\end{document}
